\documentclass{article}

\usepackage{fancyhdr}
\usepackage{extramarks}
\usepackage{amsmath}
\usepackage{amsthm}
\usepackage{amsfonts}
\usepackage{enumitem}
\usepackage{tikz}
\usepackage[plain]{algorithm}
\usepackage{algpseudocode}

\newcommand{\I}{\mathcal{I}}
\newcommand{\F}{\mathbb{F}}
\usetikzlibrary{automata,positioning}

%
% Basic Document Settings
%

\topmargin=-0.45in
\evensidemargin=0in
\oddsidemargin=0in
\textwidth=6.5in
\textheight=9.0in
\headsep=0.25in

\linespread{1.1}

\pagestyle{fancy}
\lhead{\hmwkAuthorName}
\chead{\hmwkClass\ (\hmwkClassInstructor\ \hmwkClassTime): \hmwkTitle}
\rhead{\firstxmark}
\lfoot{\lastxmark}
\cfoot{\thepage}

\renewcommand\headrulewidth{0.4pt}
\renewcommand\footrulewidth{0.4pt}

\setlength\parindent{0pt}

%
% Create Problem Sections
%


\newcommand{\enterProblemHeader}[1]{
    \nobreak\extramarks{}{}\nobreak{}
    \nobreak\extramarks{}\nobreak{}
}

\newcommand{\exitProblemHeader}[1]{
    \nobreak\extramarks{}\nobreak{}
    \stepcounter{#1}
    \nobreak\extramarks{}\nobreak{}
  }

  
\newcommand{\enterExerciseHeader}[1]{
    \nobreak\extramarks{}{}\nobreak{}
    \nobreak\extramarks{}\nobreak{}
}

\newcommand{\exitExerciseHeader}[1]{
    \nobreak\extramarks{}\nobreak{}
    \stepcounter{#1}
    \nobreak\extramarks{}\nobreak{}
}

\setcounter{secnumdepth}{0}
\newcounter{partCounter}
\newcounter{homeworkProblemCounter}
\setcounter{homeworkProblemCounter}{1}

\newcounter{homeworkExerciseCounter}
\setcounter{homeworkExerciseCounter}{1}
\nobreak\extramarks{}\nobreak{}
\nobreak\extramarks{}\nobreak{}

%
% Homework Problem Environment
%
% This environment takes an optional argument. When given, it will adjust the
% problem counter. This is useful for when the problems given for your
% assignment aren't sequential. See the last 3 problems of this template for an
% example.
%
\newenvironment{homeworkProblem}[1][-1]{
    \ifnum#1>0
        \setcounter{homeworkProblemCounter}{#1}
    \fi
    \section{Problem \arabic{homeworkProblemCounter}}
    \setcounter{partCounter}{1}
    \enterProblemHeader{homeworkProblemCounter}
}{
    \exitProblemHeader{homeworkProblemCounter}
}

  \newenvironment{homeworkExercise}[1][-1]{
    \ifnum#1>0
        \setcounter{homeworkExerciseCounter}{#1}
    \fi
    \section{Exercise \arabic{homeworkExerciseCounter}}
    \setcounter{partCounter}{1}
    \enterExerciseHeader{homeworkExerciseCounter}
}{
    \exitExerciseHeader{homeworkExerciseCounter}
}

%
% Homework Details
%   - Title
%   - Due date
%   - Class
%   - Section/Time
%   - Instructor
%   - Author
%

\newcommand{\hmwkTitle}{Homework\ \#2}
\newcommand{\hmwkAssignedDate}{January 24, 2017}
\newcommand{\hmwkDueDate}{February 5, 2017}
\newcommand{\hmwkClass}{Design and Analysis of Algorithms, CS 6550}
\newcommand{\hmwkClassTime}{}%Section A}
\newcommand{\hmwkClassInstructor}{Professor Jamie Morgenstern}
\newcommand{\hmwkAuthorName}{}%\textbf{Josh Davis} \and \textbf{Davis Josh}}

%
% Title Page
%

\title{
    \vspace{2in}
    \textmd{\textbf{\hmwkClass:\ \hmwkTitle}}\\
    \normalsize\vspace{0.1in}\small{Due\ on\ \hmwkDueDate\ at 3:10pm}\\
    \vspace{0.1in}\large{\textit{\hmwkClassInstructor\ \hmwkClassTime}}
    \vspace{3in}
}

\author{\hmwkAuthorName}
\date{}

\renewcommand{\part}[1]{\textbf{\large Part \Alph{partCounter}}\stepcounter{partCounter}\\}

%
% Various Helper Commands
%

% Useful for algorithms
\newcommand{\alg}[1]{\textsc{\bfseries \footnotesize #1}}

% For derivatives
\newcommand{\deriv}[1]{\frac{\mathrm{d}}{\mathrm{d}x} (#1)}

% For partial derivatives
\newcommand{\pderiv}[2]{\frac{\partial}{\partial #1} (#2)}

% Integral dx
\newcommand{\dx}{\mathrm{d}x}

% Alias for the Solution section header
\newcommand{\solution}{\textbf{\large Solution}}

% Probability commands: Expectation, Variance, Covariance, Bias
\newcommand{\E}{\mathrm{E}}
\newcommand{\Var}{\mathrm{Var}}
\newcommand{\Cov}{\mathrm{Cov}}
\newcommand{\Bias}{\mathrm{Bias}}

\begin{document}

% \maketitle

\pagebreak

{\bf Homework Out: \hmwkAssignedDate}\\
{\bf Due Date:  \hmwkDueDate}\\

The HW contains some exercises (fairly simple problems to check you
are on board with the concepts; don’t submit your solutions), and
problems (for which you should submit your solutions, and which will
be graded). Some problems have sub-parts that are exercises.  For this
problem set, it’s OK to work with others. (Groups of 2, maybe 3 max.)
That being said, please think about the problems yourself before
talking to others. Please cite all sources you use, and people you
work with.  Submissions due at beginning of class on the due date.
Please email \emph{LaTeXed} solutions to Tao at uthaipon3@gatech.edu.


\begin{homeworkExercise}{Exercises about MSTs.}

  \begin{enumerate}[label=\alph*.]
  \item Boruvka reduces the number of nodes by a constant factor in
    each round, but what about the number of edges? Show a graph with
    $n$ nodes and $m$ edges where the number of edges in $G_i$ remains
    $\Omega(m)$ for $\Omega(\log n)$ rounds, even after cleaning up.
  \item For a graph $G$, consider two edge-weight functions $w_1$ and
    $w_2$ such that
    %
    \[w_1(e) \leq w_1(e') \Leftrightarrow ⇐⇒ w_2(e) \leq w_2(e')\]
    %
    for all edges $e, e'\in E$. Show that $T$ is an MST wrt $w_1$ iff
    it is an MST wrt $w_2$. (In other words, only the sorted order of
    the edges matters for the MST.)
  \item Suppose graph $G$ has integer weights in the range
    $\{1, . . . , W\}$, where $W \geq 2$. Let $G_i$ be the edges of
    weight at most $i$, and $\kappa_i$ be the number of components in
    $G_i$. Then show that the MST in $G$ has weight exactly $n − W + \sum_{i= i}^{W−1}
    \kappa_i$ .
  \item Show that running (the contraction version of) Boruvka (where
    we contract components and clean up after every round) on a planar
    graph runs in linear time $O(m)$. Hint: a simple planar graph on n
    vertices has at most $3n − 6$ edges.
  \item Show that running $\log \log n$ rounds of Boruvka, followed by
    the Fibonacci heaps implementation of Prim’s, gives an
    $O(m \log \log n)$ time algorithm for MSTs.
    \end{enumerate}
  \end{homeworkExercise}

  \begin{homeworkExercise}{Exercises about arborescences (directed spanning trees).}
    \begin{enumerate}[label=\alph*.]
    \item In the rooted min-cost arborescence problem, given a digraph
      $G = (V, A)$ with edge weights $w_e$ and a root $r \in V$ , we
      want to find the min-cost $r$-arborescence. In the unrooted
      min-cost arborescence problem, given a digraph $G = (V, A)$ with
      edge weights $w_e$, we want to find the min-cost arborescence in
      the graph, regardless of which vertex it is rooted at.

      Show these problems are equivalent. I.e., given an algorithm for
      one problem, show how to solve the other problem. Your (simple)
      reductions should take linear time.

    \item Give a deterministic linear-time algorithm for finding a
      min-cost $r$-arborescence in a DAG (if one exists).

    \item For a digraph $G$, consider two arc-weight functions $w_1$
      and $w_2$ such that
      %
      \[w_1(a) \leq w_1(a') \iff w_2(a) \leq w_2(a')\]
      %
      for all arcs $a, a' \in E$. Give an example to show that the
      min-cost arborescences w.r.t. $w_1$ and $w_2$ may be different.
    \item Show an instance where $a$ is the heaviest arc in $G$, and
      there exist $r$-arborescences that do not use $a$ but the
      cheapest $r$-arborescence uses $a$. (This is not possible for
      MSTs.)
    \end{enumerate}
    
    \end{homeworkExercise}

    \begin{homeworkProblem}{(Yao’s $O(m \log \log n)$-time MST
        Algorithm)}. The idea behind this algorithm: in Boruvka’s
      algorithm we scan all the edges in the graph in each pass, and
      we should avoid this repetition.  Assume that $G$ is a connected
      simple graph, and edge weights are distinct.
      \begin{enumerate}[label=\alph*.]
      \item Suppose for each vertex, the edges adjacent to that vertex
        are stored in non-decreasing order of weights. Show a slight
        variant of Boruvka’s algorithm that runs in time
        $O(m + n \log n)$ time.
      \item ($k$-partial sorting) Given a parameter $k$ and a list of
        $N$ numbers, give an $O(N \log k)$-time algorithm that partitions
        this list into $k$ groups $g_1, g_2, . . . , g_k$ of size at most
        $\lceil N/k\rceil$ each such that all elements in $g_i$ are smaller than those
        in $g_{i+1}$ for each $i$.
      \item Adapt your algorithm from the first part above to handle
        the case where the edges adjacent to each vertex are not
        completely sorted but only $k$-partially-sorted. Ideally, your
        run-time should be $O(\frac{ m}{ k} \log n + n \log n)$.
      \item Use the two parts above (setting $k = \log
        n)$, preceded by some additional rounds of Boruvka, to give an
        $O(m \log \log n)$-time MST algorithm.
        \end{enumerate}
   
      \end{homeworkProblem}

      
      \begin{homeworkProblem}{(Fun with Boruvka) }
        In this problem, we will tie up some loose ends from lecture:
        \begin{enumerate}[label=\alph*.]
        \item (Linear-time Cleanup) Show how to implement the “shrink”
          algorithm in $O(m)$ time.  This algorithm takes as input a
          graph $G = (V, E)$ with some of the edges colored blue, and
          outputs a new graph $G' = (V' , E' )$ in which each vertex
          $v_C\in V'$ corresponds to a blue connected component $C$ in
          $G$, there are no self-loops, and there is a (single) edge $e_C
          = (v_C, v'_C)$ if there exists an edge between the
          corresponding components $C, C'$ in $G$. Moreover, the weight of
          edge $(v_C, v'_C) = \min_{\substack{\{x,y\}\in E:\\x\in C\\y\in C'}} w_{x,y}$.

        \item (Tighter Boruvka) Consider the naive implementation of
          Boruvka’s algorithm, which for each node in the current
          graph picks the cheapest edge incident to the node, and then
          shrinks the graph (as in the previous part). We saw in
          Lecture \#1 that since we spend $O(m)$ time per round, and
          there are $\leq \log_2 n$ rounds, the algorithm takes
          $O(m \log n)$ time. Show a tighter bound of
          $\min\{O(n^2), O(m(1 + \log \frac{n^2}{m}))\}$.
        \item (Boruvka on a tree) Suppose $T = (V, E)$ is a tree on
          $n$ vertices, and we run Boruvka’s algorithm on $T$. (Let
          $V_1 = V$ be the original vertices at the beginning of round
          1, $V_i$ be the vertices in round $i$, and say there $L$
          rounds so that $|V_{L+1}| = 1$). We build a tree $T'$ as
          follows: the vertices are the union of all the $V_i$. There
          is an edge from $v \in V_i$ to $w \in V_{i+1}$ if the vertex
          $v$ belongs to a component in round $i$ that is contracted
          to form $w \in V_{i+1}$; the weight of this edge $(v, w)$ is
          the min-weight edge out of $v$ in round $i$. (Note that all
          vertices in $V$ are now leaves in $T'$ .)

          For nodes $u, v$ in a
          tree $T$, let
          $\textrm{maxwt}_T (u, v)$ be the maximum weight of an edge
          on the (unique) path between $u, v$ in the tree $T$. Show that
          for all $u, v \in V$
\[\textrm{maxwt}_T (u, v) = \textrm{maxwt}_{T'} (u, v).\]
\emph{ This transformation is useful in the MST verification algorithm of
Komlos, and others.}
        \end{enumerate}
        
      \end{homeworkProblem}

      
      \begin{homeworkProblem}{(Matroids and the Greedy Algorithm)}
        Given a finite universe $U$ of elements, a collection of
        subsets $\I \subseteq 2^U$ is subset-closed if $B \in \I $ and
        $A \subseteq B \subseteq U$ implies $A \in \I$.  A set system
        $(U, \I)$ is called an independence system if {\bf
          (I1)}$\emptyset\in\I$ and {\bf (I2)} $\I$ is subset-closed. (We
        call a set $A$ independent if $A \in \I$.)

        An matroid is an independence system that satisfies the
        additional property {\bf (I3)} that if $A, B \in \I$ and $|A| < |B|$ there
        is an element $x\in  B \ A$ such that $A \cup \{x\} \in \I$.

        Some more notation: a maximal independent set is called a
        base. A cycle (or circuit) is a minimal dependent set. A cut
        is a minimal set that intersects every base. The rank of a
        matroid is the cardinality of any maximal (and hence maximum)
        independent set in the matroid.

        \begin{enumerate}[label=\alph*.]
        \item (Don’t hand in) Show that the following set systems are matroids:
          \begin{enumerate}
          \item $\I$ is the collection of all subsets of $U$ with cardinality at most $r$.
          \item $U$ is the disjoint union of sets
            $U_1, U_2, \ldots, U_\ell$ (denoted
            $U = U_1 \biguplus U_2 \biguplus \ldots \biguplus
            U_\ell$), and \[\I = \{A \subseteq U \mid  |A \cap U_i | \leq r_i \forall i \in [\ell]\}\]
          \item $U$ is the set of edges of graph $G = (V, E)$ and $\I$
            contains all acyclic subsets of edges (i.e., forests),
          \item $U = \F^r$ for some field $\F$, and $\I$ contains
            every set of vectors in $U$ that are linearly independent
            over the field $\F$.

            Show that the ranks of these matroids is
            $r,\sum_i r_i, |V| - 1$, and $r$ respectively.
          \end{enumerate}
        \item (Do not hand in) Verify that the following independence
          systems do not form matroids in general (by giving
          counterexamples):
          \begin{enumerate}
          \item $U$ is the set of edges of an undirected bipartite
            graph $G = (V, E)$, and $\I$ contains all matchings in $G$
            (i.e., subsets of edges such that no two edges share a
            common endpoint).
          \item $U$ is the set of arcs of a directed graph
            $G = (V, A)$ and $\I$ contains all subsets of arcs that do
            not contain a directed cycle.
          \end{enumerate}
        \item  (Do not hand in) Show that the cardinality of any two bases in a matroid is the same.
        \item Given $M = (U, \I)$, define the dual matroid
          $M^∗ = (U, \I^*)$ where the bases of $\I^*$ are the
          complements of bases in $M$ (and hence the independent sets
          of $\I^*$ are subsets of these complements). Show that $M^∗$
          is also a matroid.

          Given element weights $w : U \to \mathbb{R}$, the min-weight base
          problem seeks to find a base $B \in \I$ of minimum total
          weight $w(B) := \sum_{e\in B}w(e)$. (You may assume distinct
          edge weights for the following parts.)
          \begin{enumerate}
          \item (Do not hand in) Show that there is a unique min-weight base.
          \item (Do not hand in) Show that repeatedly applying the cut
            rule (pick a cut and color the lightest edge in the cut
            blue) and the cycle rule (pick any cycle and color the
            heaviest edge on the cycle red) to any matroid colors all
            the edges red/blue, and the resulting blue edges form a
            min-weight base.
          \item Consider the natural generalization of Kruskal’s
            algorithm: let $S \leftarrow\emptyset $, sort elements by
              weight and consider them in non-decreasing order, when
              considering an element $e$ add it to the $S$ if $S \cup \{e\}$ is
              independent. Show that this algorithm solves the
              min-weight base problem. (You may assume an oracle that
              given $S \subseteq U$ answers the membership query “$S \in \I$?” in
              constant time.)
            \item Show that for an independence system $(U, \I)$, if
              the greedy algorithm correctly solves the min-weight
              base problem for all settings of weights, then the
              independence system is a matroid. Hence the greedy
              algorithm characterizes matroids.
          \end{enumerate}
        \end{enumerate}
      \end{homeworkProblem}

      \begin{homeworkProblem}{(Clustering on a Tree.)}
        For a tree $T = (V, E)$ with positive edge lengths, define
        $d)_T (u, v)$ to be the shortest-path distance between $u$ and
        $v$ in $T$ according to these edge-lengths.  Given tree
        $T = (V, E)$ with $n$ nodes, and an integer
        $k\in\mathbb{Z}\geq 0$, you want to pick a set $C$ of $k$
        “centers” from $T$ such that
        \[ \sum_{v\in V}(\textrm{distance of $v$ to the closest
            center in $C$})\]
        %
        is minimized.\footnote{ If we define
          $d_T (v, C) := min_{c\in C} d_T (v, c)$, the goal is to find
          $C$ with $|C| = k$ to minimize $\sum_v  d_T (v, C)$.} Give a
        dynamic-programming solution that runs in time poly$(k, n)$.
        
        \emph{(Hint: it may help to consider the special case where each
        node in $T$ has degree 3 or less.  Also, this problem does not
        use any ideas we’ve been talking about in the class thus far.)}
       \end{homeworkProblem}
   
\end{document}
